\documentclass[10pt]{article}

\usepackage[utf8]{inputenc}

\usepackage[T1]{fontenc}

\usepackage{amsmath}

\usepackage{amsfonts}

\usepackage{amssymb}

\usepackage[version=4]{mhchem}

\usepackage{stmaryrd}

\usepackage{mathrsfs}

\usepackage{graphicx}

\usepackage[export]{adjustbox}

\graphicspath{ {./images/} }



\begin{document}

где интегрирование ведется в пределах



\[

\nsucc \leqslant H_{N}(p, q) \leqslant \mathscr{E}+\Delta \mathscr{E}

\]



Появление $h$ связано с тем, что классич. Г. р. рассматривается как предельный случай квантового (см. Квазиклассическое приближение).



Микроканонич. Г. р. не чувствительно к выбору величи-



\begin{center}

\includegraphics[max width=\textwidth]{2023_07_08_b2ff8eefc813d5310b00g-1}

\end{center}



\[

f(p, q)=A \delta\left(H_{N}(p, q)-\check{\varepsilon}\right)

\]



где $\delta$ - дельта-функция, $A-$ постоянная, определяемая из условия нормировки. Распределение вероятностей, соответствующее однородному распределению внутри слоя,микроканонич. Г. р. по энергетич. поверхности



\[

\sigma\left(H_{N}\right)=\operatorname{const}\left|\operatorname{grad} H_{N}\right|^{-1}

\]



константа определяется из условия нормировки



\[

\int \sigma\left(H_{N}\right) d \Gamma=1

\]



$\operatorname{grad} H_{N}$ берется в фазовом пространстве.



Статистич. вес $W(*, N, V)$ определяет энтропию системы $S$ как функцию \&, $N, V$ :



\[

S=R \ln W(\mathscr{\varepsilon}, N, V)

\]



а следовательно, и все другие термодинамич. функции системы, определяемой заданной формой гамильтониана.



Каноническое распределение Гиббса - распределение вероятностей состояний статистич. ансамбля систем, к-рые находятся в тепловом равновесии со средой (термостатом) и могут обмениваться с ней энергией при постоянном объеме и постоянном числе частиц, соответствует каноническому ансамблю Гиббса:



\[

f(p, q)=Z^{-1} \exp \left(-H_{N}(p, q) / k T\right),

\]



$Z=\mathrm{Z}(T, V, N)-$ статистич. интеграл, определяемый из условия нормировки $f$.



\[

Z=\int \exp \left[-H_{N}(p, q) / k T\right] d \Gamma_{N},

\]



где интегрирование ведется по фазовому пространству всех $N$ частиц. Канонич. Г. р. можно получить, если рассматривать совокупность данной системы и термостата как одну замкнутую изолированную систему и применить к ней микроканонич. Г. р. Тогда малая подсистема, функцию распределения к-рой можно найти интегрированием по фазовым переменным термостата, описывается канонич. Г.р. (те о рем а Гиббс а).



В квантовой статистич. физике статистич. ансамбль характеризуется распределением вероятностей $w_{i}$ Квантовых состояний с энергией $\varepsilon_{i}$. Канонич. Г. р. для квантовой системы имеет вид



\[

w_{i}=Z^{-1} \exp \left(-\mathscr{E}_{i} / k T\right)

\]



где $Z$ - статистич. сумма, определяемая из условия нормировки



\[

\sum_{i} w_{i}=1

\]



равна



\[

Z=\sum_{i} \exp \left(-\varkappa_{i} / k T\right)

\]



где суммирование ведется по всем квантовым состояниям допустимой симметрии. В квантовом случае канонич. Г. р. можно представить с помощью статистич. оператора (матрицы плотности)



\[

\hat{\rho}=Z^{-1} \exp (-H / k T)

\]



где $H$ - гамильтониан системы. Такая форма удобна для приложений с использованием представления вторичного квантования для гамильтониана. Канонич. Г. р. как в классическом, так и в квантовом случаях определяет с в о 6 о дную энергию (энергию Гельмгольца) в переменных $T, V, N$, равную



\[

F=-k T \ln Z

\]



где $Z$ - статистич. интеграл или статистич. сумма. Канонич. Г. р. соответствует максимуму информационной энтропии при заданной средней энергии и при сохранении нормировки.



Большое каноническое распределение Гиббса - распределение вероятностей состояний статистич. ансамбля систем, к-рые находятся в тепловом и материальном равновесии со средой (термостатом и резервуаром частиц) и могут обмениваться с ними при постоянном объеме $V$, соответствует большому каноническому ансамблюГиббса



\[

f_{N}(p, q)=Z^{-1} \exp \left(-\left(H_{N}(p, q)-\mu N\right) / k T\right)

\]



где $\mu-$ химич. потенциал, $Z=Z(T, \mu, V)-$ статистич. интеграл, определяемый из условия нормировки



\[

\sum_{N \geqslant 0} \int f_{N} d \Gamma_{N}=1

\]



равен



\[

Z=\sum_{N \geqslant 0} \int \exp (-(H(p, q)-\mu N) / k T) d \Gamma_{N}

\]



Болышое канонич. Г. р. можно вывести, если рассматривать совокупность данной системы вместе с термостатом и резервуаром частиц как большую изолированную и замкнутую систему и применить к ней микроканонич. Г. р. Тогда малая подсистема описывается болышим канонич. Г.р., к-рое можно найти интегрированием по фазовым переменным термостата и резервуара частиц и суммированием по числу частиц (т е о р е ма Ги б б с а).



В квантовой статистике для большого канонич. Г. р.



\[

w_{i N}=Z^{-1} \exp \left(-\left(\varkappa_{i}-\mu N\right) / k T\right)

\]



где условие нормировки



\[

\sum_{i, N} w_{i n}=1

\]



И



\[

Z=\sum_{i, N \geqslant 0} \exp \left(-\left(\varkappa_{i}-\mu N\right) / k T\right)

\]



а статистич. оператор



\[

\hat{\rho}=Z^{-1} \exp (-(H-\mu N) / k T) .

\]



Большое канонич. Г. р. как в классическом, так и в квантовом случае определяет термодинамич. потенциал



\[

\Omega=-k T \ln Z

\]



в переменных $T, \mu, V$, к-рый наиболее удобен для приложений, так как содержит лишь один экстенсивный параметр $V$. Большое канонич. Г. р. можно получить, если рассматривать совокупность данной системы, являющейся также резервуаром частиц, как одну замкнутую изолированную систему (т е о е ма Гиббса для большого канонич. Г.р.).



Большое канонич. распределение соответствует максимуму информационной энтропии при заданных средней энергии и среднего числа частиц, а канонич. распределение при заданной средней энтропии. В обоих случаях нормировка предполагается фиксированной (теоремы об экстремальных свойствах канонич. распределений, доказанные Дж. Гиббсом, 1901).



Термодинамическая эквивалентность распределений Гиббca. Все Г. р. вдали от областей сосуществования фаз и критич. точек, вследствие малости относительных флуктуаций, пропорциональных $N^{-1 / 2}$, термодинамически эквивалентны, то есть приводят к одинаковым термодинамич. функциям, после приведения их к одинаковым переменным с помощью преобразования Лежандра. В области фазовых переходов, вообще говоря, нет эквивалентности Г.р. относительно вычисления корреляционных функщий и функций распределения флуктуаций.



В неравновесной статистич. механике возможно также использование идей статистич. ансамблей Гиббса, напр. с помощью отбора запаздывающих решений уравнения Лиувилля через граничные условия.



ГИББСА 131





\end{document}